% !PRINTED 560
\weberpage{560}

%<<CONT
\[
\begin{aligned}
\frac{\Gamma(2s)}{[-2\pi i(x+\omega y)]^{2s}} &= \int_0^\infty e^{2\pi i(x+y\omega)\xi}\xi^{2s-1}\,d\xi,\\
&= \int_0^1 e^{2\pi i x\xi}\sum_{0,\infty}^n e^{2\pi i y\omega(\xi+n)}(\xi+n)^{2s-1}\,d\xi.
\end{aligned}
\]
If one now sums over $x$ and applies the \gesperrt{Fourier} series:
\[
\sum^x \int_0^1 e^{2\pi i x\xi} f(\xi)\,d\xi = \tfrac{1}{2}[f(0)+f(1)],
\]
then it follows:
\[
\frac{\Gamma(2s)}{(-2\pi i)^{2s}} \sum_{-\infty,\infty}^x \frac{1}{(x+\omega y)^{2s}} = \sum_{1,\infty}^n e^{2\pi i y\omega n}\,n^{2s-1},
\]
and by summation over $y$:
\[
\begin{aligned}
\frac{\Gamma(2s)}{(-2\pi i)^{2s}} \sum_{1,\infty}^y \sum_{-\infty,\infty}^x \frac{1}{(x+\omega y)^{2s}} &= \sum_{1,\infty}^n \frac{e^{2\pi i\omega n}}{1-e^{2\pi i\omega n}}\,n^{2s-1}\\
&= -\frac{1}{2\pi \imath}\sum_{1,\infty}^n \frac{d\log(1-e^{2\pi i\omega n})}{d\omega}\,n^{2s-2}.
\end{aligned}
\]
Accordingly we obtain, by passing to the limit and having regard to the definition of the function $\eta(\omega)$ [§ 24, (8)]:
\[
\mathop{\mathrm{Lim}}\limits_{s=1}\sum_{1,\infty}^y \sum^x \frac{1}{(x+\omega y)^{2s}} = -2\pi i\,\frac{d\log\eta(\omega)}{d\omega} - \frac{\pi^2}{6}.
\]
Furthermore we have
\[
\mathrm{Lim}\sum \frac{1}{x^{2s}} = \frac{\pi^2}{6},
\]
and accordingly from (1) and (2) it follows:
\weq{3}{\mathop{\mathrm{Lim}}\limits_{s=1}\sum^{x,y} \frac{1}{(x+\omega y)^{2s}} = -4\pi i\,\frac{d\log\eta(\omega)}{d\omega}}
[to be summed as indicated for (1)]. If one wished here to pass to the limit under the summation sign, one would obtain:
\weq{4}{-4\pi i\,\frac{d\log\eta(\omega)}{d\omega} = \sum^{x,y} \frac{1}{(x+\omega y)^2},}
but then one would have a sum that is only conditionally convergent, and a precise manner of passing to the limit would have to be laid down. One would have to assume in the $xy$-plane a closed curve within which the points $x, y$ lie, and then let this curve move out to infinity. The value of the sum will depend on the nature of this curve. Perhaps the ellipse $|x+\omega y| = n$ presents itself for this. In itself formula (4) has no meaning.

\begin{center}\rule{3cm}{0.4pt}\end{center}
% [sic] Print reads'und demnach folgt(1) und(2)' without'aus'; German reading retained, English'from' makes the evident relation explicit.
